\documentclass[11pt,a4paper]{article}

% --- Packages ---
\usepackage[utf8]{inputenc}
\usepackage[T1]{fontenc}
\usepackage{geometry}
\geometry{margin=2cm}
\usepackage{array}
\usepackage{fancyhdr}
\usepackage{titlesec}
\usepackage{tcolorbox}
\tcbuselibrary{listings,skins,breakable}
\usepackage{amsmath,amssymb}
\usepackage{float}
\usepackage{xurl}
\usepackage[hidelinks]{hyperref}
\usepackage{graphicx}
\usepackage{subcaption}
\usepackage[english]{babel}
\usepackage{mwe}

\lstdefinestyle{pythonstyle}{
  language=Python,
  backgroundcolor=\color{gray!5},
  basicstyle=\ttfamily\small,
  keywordstyle=\color{blue}\bfseries,
  commentstyle=\color{gray}\itshape,
  stringstyle=\color{teal},
  numberstyle=\tiny\color{gray},
  numbers=left,
  stepnumber=1,
  numbersep=8pt,
  showspaces=false,
  showstringspaces=false,
  showtabs=false,
  tabsize=4,
  breaklines=true,
  breakatwhitespace=true,
  frame=single,
  rulecolor=\color{gray!40},
  captionpos=b,
  keepspaces=true,
  morekeywords={ceil,cbrt,round,split,factorial,pi},
}

\hypersetup{
  pdftitle={I Hate Mathematics},
  pdfauthor={Ludwig Alvarado},
  pdfsubject={Why someone hate mathematics? They're the most beautiful thing in the world, it's like hating your mother...},
  pdfkeywords={data structures and algorithms},
  pdfcreator={GNU Emacs desde Arch (btw) GNU/Linux y LaTeX con hyperref},
  pdfproducer={pdflatex}
}


% --- Formatting ---
\pagestyle{fancy}
\fancyhf{}
\setlength{\headheight}{30pt}
\setlength{\footskip}{20pt}
\setlength{\parindent}{0pt}
\setlength{\parskip}{0.7em}

% Left header (university and exam info)
\rhead{%
  \small
  University Jorge Tadeo Lozano\\
  Data Structures - Group 2\\
  Programming Test 
}

% Right header (date, professor, monitor)
\lhead{%
  \small
  \textbf{Date}: \today \\
  \textbf{Professor:} José Alejandro Franco Calderón \\
  \textbf{Assistant:} Ludwig Alvarado Becerra
}

\cfoot{\thepage}

\titleformat{\section}{\large\bfseries}{\thesection.}{0.5em}{}

% --- Environment for samples ---
\newenvironment{sample}[2]{
  \begin{tcolorbox}[title=Prueba \##1, colback=white, colframe=black, sharp corners,
    fonttitle=\bfseries, breakable]
    \begin{tabular}{|p{0.45\linewidth}|p{0.45\linewidth}|}
      \hline
      \textbf{Input} & \textbf{Output} \\
      \hline
      }{\\ \hline
    \end{tabular}
  \end{tcolorbox}
}

% --- Document ---
\begin{document}

\begin{center}
  {\LARGE \bf I Hate Mathematics}\\[1ex]
  Author: Ludwig Alvarado
\end{center}

Hi, my name is Vectra, I am a normal university student at the University of Tears, Anxiety, Deadlines, Exams and Overthinking (UTADEO) who really really hates math. My degree is systems engineering but all math doesn't matter there, I don't think that this bunch of math is ever going to be useful in life.

Me and my friends love to \textit{parchar} in ``La Brújula''. The other day, a random guy came to us shouting a bunch of mathematical formulas that I didn't understand. In his madness he said that all of those formulas lead to the secret of the universe.

Although I don't like math, I do feel curious about the universe... can you help me calculate the numeric results of the formulas that he said?

You are given several queries. Each query corresponds to a specific mathematical formula. Your task is to evaluate each query and print its result.


\begin{itemize}
  \item \textbf{ADD}$(a,b) = a + b$
  \item \textbf{MUL}$(a,b) = a \times b$
  \item \textbf{POW}$(a,b) = a^b$ (using a loop, NOT built-in power)
  \item \textbf{FACT}$(n) = n!$ (using a recursive function)
  \item \textbf{POLY}$(a,b) = a^2 + 2ab + b^2$
  \item \textbf{DIST}$(x_1,y_1,x_2,y_2) = (x_1 - x_2)^2 + (y_1 - y_2)^2$
  \item \textbf{AVG}\( (a,b,c) =  \frac{a+b+c}{3}\) 
  \item \textbf{SECRET OF UNIVERSE}\((n,a,b) =\)
        \[
\sum_{i=1}^{n}
\left(
\sum_{j=1}^{n}
\left(
\prod_{k=1}^{n}
\left(\pi \times e^{i + j +k} +  (a^2 + 2ab + b^2) + \frac{a + b + n}{3} \right)
\right)
\right)
\]
\end{itemize}


\subsection*{Program input}

The first line contains an integer \(q\) \((1 \leq q \leq 100)\), the number of queries.

Each of the next \(q\) lines contains one of the following queries:

\begin{itemize}
  \item \texttt{ADD a b}
  \item \texttt{MUL a b}
  \item \texttt{POW a b}
  \item \texttt{FACT n}
  \item \texttt{POLY a b}
  \item \texttt{DIST x1 y1 x2 y2}
  \item \texttt{AVG a b c}
  \item \texttt{UNIVERSE n a b}
\end{itemize}

All values are integers in the range:
\begin{itemize}
  \item General values: \([-10^3, 10^3]\)
  \item Factorial: \(0 \leq n \leq 10\)
  \item SECRET OF UNIVERSE: \(1 \leq n,a,b \leq 30\)
\end{itemize}

\subsection*{Program output}

For each query, print a single line with the result. Any decimal value must be rounded to two decimal places.

\subsection*{Sample test cases}

\begin{sample}{1}{}
\texttt{8}

\texttt{ADD 1 1}

\texttt{MUL 2 2}

\texttt{POW 3 3}

\texttt{POLY 2 3}

\texttt{DIST 0 0 3 4}

\texttt{AVG 3 4 8}

\texttt{FACT 5}

\texttt{UNIVERSE 2 2 1}

&
\texttt{2}

\texttt{4}

\texttt{27}

\texttt{25}

\texttt{25}

\texttt{5}

\texttt{120}

\texttt{796762.03}

\end{sample}

\subsection*{Notes}

\begin{itemize}
  \item You must build complex functions using simpler ones.
  \item Avoid using built-in shortcuts such as \texttt{**}, \lstinline[style=pythonstyle]{math.pow()}, or \lstinline[style=pythonstyle]|math.factorial()|.
  \item You can use constant such as \lstinline[style=pythonstyle]|math.e| for Euler number (\(e\)) and \lstinline[style=pythonstyle]|math.pi| for pi \(\pi\).
\end{itemize}

\subsection*{Important aspects}

\begin{enumerate}
  \item The exam will be completed in \textbf{2 stages}:
        \begin{itemize}
          \item Algorithm development on paper \textbf{(3.0 points)}.
                \begin{itemize}
                  \item You must write and justify the time and memory complexity of your algorithm.
                  \item You must identify the complexity of each function and the overall solution.
                \end{itemize}
          \item Implementation in the code editor \textbf{(2.0 points)}.
        \end{itemize}

  \item The data on each query line is provided as a string. You may use:
\begin{lstlisting}[style=pythonstyle]
data = input().split()
\end{lstlisting}

  \item This exercise evaluates the topics covered up to \textbf{functions and loops}.

  \item You must avoid lambda functions.

  \item It is not necessary to use advanced data structures or external libraries.

  \item Output format must match exactly.
\end{enumerate}

\end{document}
